\section{SDS 的问题与挑战}

尽管 SDS 取得了令人印象深刻的结果，但它也存在一些固有的问题。

\subsection{Janus Problem（多面问题）}

SDS 最著名的失败模式是 \textbf{Janus Problem}（以古罗马双面神 Janus 命名）：生成的 3D 模型从每个视角看都像正面，导致一个物体出现多张脸。

\begin{warningbox}{Janus Problem 的根本原因}
Janus Problem 的根源在于 2D 先验与 3D 一致性之间的矛盾：

\begin{enumerate}
    \item 图像扩散模型在\textbf{单视角}图像上训练，不具备 3D 一致性的概念
    \item SDS 独立优化每个视角的渲染结果，每个视角都倾向于生成该文本描述最``典型''的图像
    \item 对于人脸等有明显正面的物体，``最典型''的图像几乎总是正面视角
    \item 结果就是 3D 模型的每个面都变成了正面——产生多张脸
\end{enumerate}

本质上，这是因为 2D 扩散模型只有\textbf{局部观测}（单视角图像）的先验，而 3D 生成需要\textbf{全局一致性}的约束。
\end{warningbox}

\subsection{CFG 权重与多样性的权衡}

在使用 SDS 时，通常需要设置非常高的 Classifier-Free Guidance（CFG）权重才能获得干净的输出。

\begin{knowledgebox}{高 CFG 权重的双刃剑效应}
高 CFG 权重会将生成分布收缩到与文本描述高度匹配的狭窄区域：

\begin{itemize}
    \item \textbf{优点}：输出质量更高、与文本描述的一致性更强
    \item \textbf{缺点}：多样性大幅降低，所有生成结果趋于相似
\end{itemize}

如何在保证质量的同时维持多样性，是 SDS 的一个重要开放问题。
\end{knowledgebox}

\subsection{其他局限性}

\begin{itemize}
    \item \textbf{过饱和问题}：SDS 生成的纹理常常过于饱和、缺乏自然感
    \item \textbf{优化不稳定}：由于每次迭代的随机视角和随机噪声，优化过程可能振荡
    \item \textbf{计算成本}：每次迭代需要一次扩散模型的前向传播，整个优化过程需要数千次迭代
    \item \textbf{几何质量}：生成的 3D 几何往往不够精细，尤其是遮挡区域
\end{itemize}

\begin{warningbox}{从局部先验到全局一致：本质困难}
SDS 的所有问题都可以追溯到同一个根本困难：用\textbf{局部的 2D 观测先验}去约束\textbf{全局的 3D 结构}。这是一个病态（ill-posed）问题——无数种 3D 结构都能产生局部看起来合理的 2D 渲染。仅靠 2D 先验不足以消除这种歧义，还需要引入 3D 领域特定的先验知识。
\end{warningbox}

\subsection{本章小结}

SDS 的主要挑战包括 Janus Problem、过饱和纹理、CFG 权重调节等。这些问题的根源在于 2D 先验的局部性无法完全约束 3D 结构的全局一致性。后续工作（如 SDS 的各种改进版本 VSD、ISM 等）致力于缓解这些问题。

\newpage
